\section{Projective cancellation, weights, and ranks}\label{sec:weights}

\subsection{The four radial strata}

For a nonzero direction \([x]\in\PP^{2n-1}(\Fp)\),
\[
 \Phi_{p,n}(tx)=t^2\{2t\cC_n(x)+\cQ_{n,\rho}(x)\}.
\]
The nonzero roots are counted by four disjoint cases:
\[
\begin{array}{c|c}
\text{condition on }[x]&\text{number of nonzero }t\\ \hline
\cC_n\ne0,\ \cQ_{n,\rho}\ne0&1\\
\cC_n\ne0,\ \cQ_{n,\rho}=0&0\\
\cC_n=0,\ \cQ_{n,\rho}\ne0&0\\
\cC_n=\cQ_{n,\rho}=0&p-1.
\end{array}
\]
It follows that
\begin{equation}\label{eq:radial}
\begin{aligned}
 Z_{p,n}
 &=1+\#\{\PP^{2n-1}\setminus(S_n\cup Q_n)\}+(p-1)\#X_n\\
 &=1+P_{2n-1}-\#S_n-\#Q_n+p\#X_n,
\end{aligned}
\end{equation}
where \(P_m=1+p+\cdots+p^m\).

The quadric \(Q_n\) is split in the source rows.  Write
\[
 \#S_n=P_{2n-2}+A_{p,n},\qquad
 \#Q_n=P_{2n-2}+p^{n-1},
\]
\[
 \#X_n=P_{2n-3}-B_{p,n}.
\]
Substitution in \eqref{eq:radial} cancels every projective Tate polynomial:
\begin{equation}\label{eq:z-cancel}
 Z_{p,n}=p^{2n-1}-p^{n-1}-A_{p,n}-pB_{p,n}.
\end{equation}
Therefore
\begin{equation}\label{eq:c-cancel}
 C_{p,n}=-2-\frac{2A_{p,n}}{p^{n-1}}
              -\frac{2B_{p,n}}{p^{n-2}}.
\end{equation}

The two degree-one primes above a split \(p\) give the same data.  Indeed,
\[
 x_0=\rho^2y_{2n-1},\qquad x_i=y_{i-1}\quad(1\le i\le2n-1)
\]
preserves \(\cC_n\) and sends the \(\rho\)-closing quadric to the
\(\rho^2\)-closing quadric.

\subsection{Purity and the trace identity}

Because \(S_n\) has even dimension \(2n-2\), its primitive middle trace
enters its point count with positive sign.  Because \(X_n\) has odd
dimension \(2n-3\), its middle trace enters with negative sign.  Hence
\[
 A_{p,n}=\Tr(F_p\mid H^{2n-2}_{\mathrm{prim}}(S_n)),\qquad
 B_{p,n}=\Tr(F_p\mid H^{2n-3}(X_n)).
\]
After the twists in \eqref{eq:packets},
\[
 e_{p,n}=1+\frac{A_{p,n}}{p^{n-1}},\qquad
 o_{p,n}=\frac{B_{p,n}}{p^{n-2}}
\]
are the traces of \(E_n\) and \(O_n\).  Deligne's theorem
\cite[Theorem 1.6]{Deligne1974} makes the packets pure of weights zero
and one.  Equation \eqref{eq:c-cancel} becomes
\[
 C_{p,n}=-2(e_{p,n}+o_{p,n}),
\]
proving \eqref{eq:trace-main}.  In particular,
\[
 e_{p,n}=O(1),\qquad o_{p,n}=O(p^{1/2}).
\]

\subsection{A conditional family identity}

For a smooth cubic hypersurface in \(\PP^{2n-1}\), the primitive Hodge
pieces are computed by the Jacobian ring
\[
 R_n=\CC[x_0,\ldots,x_{2n-1}]/(x_0^2,\ldots,x_{2n-1}^2),
\]
whose Hilbert series is \((1+t)^{2n}\)
\cite[§10, especially Theorem 10.8]{Griffiths1969II}.  The resulting
primitive middle Betti number is
\begin{equation}\label{eq:cubic-rank}
 b^{\mathrm{prim}}_{2n-2}(S_n)=\frac{4^n+2}{3}.
\end{equation}

For a smooth \(X_n=(2,3)\subset\PP^{2n-1}\), the normal sequence gives
\[
 c(TX_n)=\frac{(1+H)^{2n}}{(1+2H)(1+3H)},\qquad \deg X_n=6.
\]
Thus
\[
 \chi(X_n)=6[H^{2n-3}]
 \frac{(1+H)^{2n}}{(1+2H)(1+3H)}.
\]
Using
\[
 \frac1{(1+2H)(1+3H)}
 =-\frac2{1+2H}+\frac3{1+3H},
\]
two finite binomial sums give
\[
 [H^{2n-3}]
 \frac{(1+H)^{2n}}{(1+2H)(1+3H)}
 =\frac{3n+1-4^n}{9}.
\]
Consequently,
\[
 \chi(X_n)=\frac{2(3n+1-4^n)}3.
\]
Weak Lefschetz \cite[Expos\'e XIV, Corollaire 4.6]{SGA2} leaves only
the odd middle Betti number undetermined and gives
\begin{equation}\label{eq:intersection-rank}
 b_{2n-3}(X_n)=(2n-2)-\chi(X_n)
 =\frac{2(4^n-4)}3.
\end{equation}
Equations \eqref{eq:cubic-rank}--\eqref{eq:intersection-rank}, together
with the trivial line in \(E_n\), prove the rank assertions in
\Cref{thm:main}.

\begin{center}
\begin{tabular}{@{}rrrr@{}}
\toprule
\(n\)&\(\rank E_n\)&\(\rank O_n\)&total\\
\midrule
2&7&8&15\\
3&23&40&63\\
4&87&168&255\\
\bottomrule
\end{tabular}
\end{center}

The computation can be replayed symbolically for larger \(n\), but it is
not a smoothness proof for the corresponding H\'enon source members.
